%% kapa: exam=1 round=37 year=2026 period=1 form=A subject=civil_law subject_ko=민법 seq=25 number=25
%% source: Q-Net 감정평가사 기출문제 (cst00309) artlSeq=5258924 file=raw/1008/5258924/2026년 제37회 감정평가사 1차 1교시.pdf
%% conversion: kordoc 4.24.51 -> tools/gen1.py (auto); review=proofread-sonnet; flags=none
\begin{question}{25}
선의취득에 관한 설명으로 옳지 않은 것은? (다툼이 있으면 판례에 따름)
\begin{choices}
  \item 저당권은 선의취득의 대상이 되지 않는다.
  \item 민법 제249조가 규정하는 선의취득의 요건인 선의ㆍ무과실의 기준시점은 물권행위가 완성되는 때이다.
  \item 채무자 이외의 자의 소유에 속하는 동산이 경매절차에서 매각된 경우, 매수인이 선의ㆍ무과실로 매수대금을 완납하더라도 매수인은 그 동산의 소유권을 선의취득할 수 없다.
  \item 동산의 선의취득은 양도인이 무권리자라고 하는 점을 제외하고는 아무런 흠이 없는 거래행위이어야 성립한다.
  \item 점유개정의 방법으로 양도담보를 설정한 동산소유자가 다시 제3자와 양도담보설정계약을 맺고 그 동산을 점유개정으로 인도한 경우, 그 제3자는 양도담보권을 선의취득하지 못한다.
\end{choices}
\end{question}
\answer{3}
